\cleardoublepage

\chapter*{Agradecimientos}

En primer lugar, quiero agradecer este trabajo a mi familia, y de manera especial a mis padres. Gracias por vuestro apoyo constante, por vuestra paciencia infinita, por vuestra confianza incluso en los momentos en los que a mí me faltaba, y por haber estado siempre disponibles de todas las formas posibles: emocional, física, económica y, sobre todo, amorosamente. Gracias por acompañarme en tantas noches de angustia, cansancio y agobio, por ayudarme a volver a levantarme una y otra vez, y por demostrarme vuestro cariño sin reproches, tanto en los momentos buenos como en los más difíciles.

A mi padre, por ser una de mis mayores fuentes de inspiración. Por enseñarme, con su ejemplo, el valor del esfuerzo, la responsabilidad y la vocación por la ingeniería. Gran parte de mi forma de mirar este proyecto nace de haber querido aprender de él y parecerme, aunque sea un poco, a la persona y al profesional que siempre he admirado.

A mi madre, por ser calma, refugio y equilibrio. Por saber poner paz cuando todo parecía hacerse grande, por escucharme, por cuidarme y por ayudarme a ordenar no solo las ideas, sino también los miedos. Su bondad, su sensibilidad y su forma de acompañarme han sido una inspiración constante, también en la manera de escribir, de expresar y de dar sentido a este trabajo.

A los dos, gracias por los valores, la educación y el amor con los que me habéis hecho crecer. Que creáis ciegamente en mí me ha dado fuerza incluso cuando yo no la encontraba.

A mi hermana, por saber sacarme una sonrisa cuando más lo necesitaba. Por distraerme como nadie, por entenderme sin demasiadas explicaciones y por ese vínculo nuestro que no se parece a ningún otro. Estoy segura de que será una enfermera extraordinaria, porque cuidar de los demás ya forma parte de ella.

A mis abuelas, por su cariño y apoyo incondicional. Aunque no siempre entendieran todos los detalles técnicos del proyecto, siempre supieron valorar la ilusión que había detrás y acompañarme desde el amor, el orgullo y la confianza.

A mis tíos, por su cariño, su ilusión y su forma de escucharme cada vez que hablaba del proyecto, incluso cuando no era fácil explicar todo lo que había detrás. Y, de manera muy especial, a mis abuelos: a mi abuelo Andrés, a quien tuve la suerte de conocer durante muchos años y de quien sé que me cuida desde arriba, dándome fuerza y apoyo; y a mi abuelo José Luis, con quien compartí menos tiempo del que me habría gustado, pero de quien estoy segura de que habría sido un abuelo maravilloso.

A mi abuela Victoria, mi mejor amiga, aunque quizá ella no sepa hasta qué punto lo es. Gracias por cada conversación, por cada consejo de amor, de madre, de amiga, de confidente y de abuela a la vez. Hablar con ella de su vida, escucharla y compartir esos momentos me recuerda la suerte que tengo de tenerla y me hace sentir profundamente orgullosa de mis raíces. Ojalá algún día pueda parecerme un poco a ella: a su forma de querer, de cuidar, de aconsejar y de estar. Y gracias también a esa amiga que tenemos ahí arriba, como nosotras decimos, por darnos fuerza, esperanza y el regalo de seguir compartiendo tiempo juntas cuando más lo necesitábamos.

A Jorge, por estar siempre ahí, incluso cuando yo no sabía muy bien cómo estar. Gracias por acompañarme con paciencia, amor y respeto; por darme espacio cuando lo necesitaba y por acercarte cuando más falta me hacía. Por entender mis silencios, mis nervios, mis agobios y mis días de cansancio sin hacerme sentir nunca una carga. Gracias por sostenerme en los momentos difíciles, por celebrar conmigo cada pequeño avance y por recordarme, con esa calma tan tuya, que podía con esto. Su forma de quererme, constante, generosa y tranquila, ha sido uno de mis mayores refugios durante este camino.

Quiero agradecer también a los trabajadores de Agrolab Móstoles su disponibilidad, cercanía e ilusión al recibir el proyecto. Su interés, sus comentarios y su buena disposición hicieron que la prueba en un entorno real tuviera un valor especial, no solo técnico, sino también humano.

Por último, pero no menos importante, quiero expresar mi agradecimiento a D. Javier Simó, director y tutor de este Trabajo de Fin de Grado, por su ayuda, su orientación, su confianza en el proyecto y su disponibilidad durante todo el proceso. Su apoyo técnico y sus ánimos han sido fundamentales para avanzar con criterio y seguridad.

Asimismo, quiero agradecer a Dña. Ana Zazo, co-tutora de este trabajo, la oportunidad de acercar TierraViva al entorno de Agrolab Móstoles y de enriquecer el proyecto con una perspectiva más agrícola, urbana y social.

\begin{flushright}
		\emph{A mi abuela, mi mejor amiga;\\
      nuestra Toyi}\\
		\par
		Madrid, \today\\ 
		\emph{Lucía Tejedor}
\end{flushright}

\thispagestyle{empty}

